\clearpage
\fancyhead[L]{\small LAPLACE. EXPOSITION DU SYSTÈME DU MONDE.}
\SourceLeaf{505}
\begin{center}\small Göttingische gelehrte Anzeigen. Stück 119. S. 1185..1192. 23. Juli 1808.\\[3pt]\rule{.72\linewidth}{.3pt}\end{center}

\emph{Exposition du système du monde, par} M. \textsc{Laplace}, \emph{chancelier du Senat-Conservateur etc. Troisième édition, revue et augmentée par l'auteur. Chez Courcier,} 1808. 405 Seiten im Quartformat. Mit einem Bildnis des Verfassers.

Gegenstand, Aufbau, Darstellungsweise und Wert dieses in seiner Art klassischen Werks, dessen frühere Ausgaben sich in aller Händen befinden, sind zu bekannt, als dass wir uns jetzt noch damit aufhalten dürften. Eine ausführliche Besprechung seines Inhalts nach der deutschen Übersetzung der ersten Ausgabe findet sich bereits in Stück 97 und 131 dieser Blätter aus dem Jahr 1798. Die zweite Ausgabe, die 1799 erschien, hatte keine bedeutenden Änderungen erfahren. Wir beschränken uns deshalb hier auf die wichtigsten Änderungen und Zusätze, durch die sich die dritte Ausgabe von den früheren unterscheidet. Schon die größere Seitenzahl zeigt, dass die Erweiterungen erheblich sein müssen, denn die zweite Ausgabe hatte nur 351 Seiten. Die Erweiterungen bestehen teils aus einigen größeren Zusätzen, teils aus vielen kleineren. Dadurch werden die behandelten Gegenstände vollständiger erfasst und ihre Darstellung erhält einen noch engeren Zusammenhang, wird einsichtiger, fruchtbarer und interessanter. Ein besonderer Vorzug dieses Werks besteht darin, dass es von den numerischen Ergebnissen der Astronomie, soweit sie in eine allgemeine Übersicht gehören können, die neuesten und bewährtesten Bestimmungen mitteilt. In dieser Hinsicht kann es umso eher als maßgeblich gelten, als sehr viele davon auf den eigenen gründlichen Untersuchungen des Verfassers beruhen oder durch sie angeregt wurden. Wir glauben, manchem Leser einen willkommenen Dienst zu erweisen, wenn wir in dieser Besprechung die wichtigsten dieser Bestimmungen hervorheben, die seit der 2. Ausgabe weiter verbessert worden sind.

Das \emph{erste Buch}, das sich mit den \emph{scheinbaren} Bewegungen der Himmelskörper befasst, hat zwei neue Kapitel erhalten, nämlich \emph{das zehnte} über die teleskopischen Planeten Ceres, Pallas, Juno und Vesta und \emph{das elfte} über die Bewegung der Planeten um die Sonne. Das zweite Kapitel des zweiten Buchs mit derselben Überschrift ist dagegen entfallen. Tatsächlich treten die Gründe für die wahre Weltordnung deutlicher hervor, wenn vorher gezeigt worden ist, dass uns die Bewegungen der Planeten genau so erscheinen, als fänden sie auf Epizykeln statt, deren Mittelpunkt stets mit der Sonne zusammenfiele und sich mit ihr um die Erde bewegte. Das lässt sich bei Venus und Merkur aus ihren Phasen zusammen mit den aufeinanderfolgenden Änderungen ihrer scheinbaren Durchmesser beweisen, bei Jupiter und Saturn aus den Verfinsterungen der Trabanten und dem wiederholten Verschwinden des Rings. Diese Folgerungen aus den Tatsachen gehören also durchaus in den ersten Abschnitt, da sie von der Wahl des Weltsystems unabhängig sind. Vor der Erfindung der Fernrohre waren diese Tatsachen freilich unbekannt. \textsc{Ptolemäus} konnte daher in seinem System bei jedem Planeten nur das Verhältnis der Radien der Kreise und Epizykel angeben und musste ihre absoluten Größen unbestimmt lassen.

Von den verbesserten numerischen Angaben in diesem Buch nennen wir die folgenden und beziehen uns dabei auf die Seitenzahlen der zweiten Ausgabe. Dass überall die neuen Dezimaleinteilungen zugrunde liegen, braucht kaum erwähnt zu werden. Auf S. 6 beträgt die Schiefe der Ekliptik für 1801 $26^{\circ}\,07315$, ihre säkulare Abnahme (S. 11) für das gegenwärtige Jahrhundert $160''85$. Auf S. 16 beträgt der siderische Tag 0.997269672 des mittleren Sonnentags. Auf S. 17 dauert das tropische Jahr jetzt 365.242264 Tage. Auf S. 20 dauert der siderische Umlauf des Mondes 27.3216610716 Tage zu Beginn des 19. Jahrhunderts, der siderische Umlauf der Erdnähe 3232.58075
\SourceLeaf{506} Tage, der des Knotens 6793.42118 Tage. Für die Epoche von 1801 (Mitternacht vor dem 1. Januar in Paris) beträgt die mittlere Länge des Mondes $124^{\circ}01299$, die der Erdnähe $295^{\circ}66824$ und die des aufsteigenden Knotens $17^{\circ}6933$. Ferner beträgt die Exzentrizität der Mondbahn 0.0548553, die Neigung $5^{\circ}7222$ und der synodische Umlauf des Mondes 29.58058817896 Tage. Die neuen Bestimmungen zu den Planeten stellen wir weiter unten zusammen. Auf S. 51 beträgt die große Achse der Nutationsellipse $59''56$, die kleine in Teilen des Parallelkreises $111''30$. Auf S. 55 beträgt der nördliche Grad bei der mittleren Breite $73^{\circ}7$ nach den französischen Astronomen 100696 Meter, nach der neuen schwedischen Gradmessung aber nur 100316.1 Meter. Vielleicht würde sich bei einer erneuten Messung auch der von \textsc{Lacaille} am Kap gemessene Grad nicht unbeträchtlich ändern. Auf ihn scheint man wohl etwas zu viel Gewicht gelegt zu haben, um die Unregelmäßigkeit der Erdgestalt zu beweisen. Der Darstellung der wichtigsten Methoden zur Bestimmung der Länge hat \textsc{Laplace} den Wunsch beigefügt, dass sich alle Nationen darauf einigen möchten, die Länge von einem besonders auffälligen natürlichen Punkt aus zu zählen, statt von der wichtigsten Sternwarte des jeweiligen Landes. Dafür würde sich der Gipfel von Teneriffa besonders gut eignen. Auf S. 81 beträgt das Verhältnis des spezifischen Gewichts der atmosphärischen Luft zu dem des Quecksilbers beim Barometerstand 0.76 und bei der Temperatur schmelzenden Eises 1 zu 10477.9. Der ganze Abschnitt über Höhenmessungen mit dem Barometer und über die Lichtbrechung wird ausführlicher behandelt.

Im \emph{zweiten Buch} über die wahren Bewegungen der Himmelskörper finden wir nur wenige Änderungen. Im VI. (jetzt V.) Kapitel ist die Schätzung der Wahrscheinlichkeit entfallen, dass zwei beobachtete Kometen mit nahezu gleichen Elementen ein und derselbe sind. Stattdessen werden einige Vermutungen über die Wirkung der Sonnenwärme auf die Kometen gewagt. Völlig umgearbeitet ist dagegen die Tafel mit den Bestimmungsgrößen sämtlicher Planeten, die wir wegen ihrer Bedeutung hier vollständig wiedergeben. Dass \textsc{Laplace} die Epoche jetzt von Mitternacht zwischen dem 31. Dezember und dem 1. Januar an zählt, wurde bereits erwähnt.

\begin{center} \Needspace{140pt}
{\small\setlength{\tabcolsep}{6pt}
\begin{tabular}{l|r|r}
 & \shortstack{Mittlere Länge 1801\\Pariser Meridian} & \shortstack{Siderischer Umlauf\\in Tagen} \\
\hline
Merkur & $182^{\circ}15647$ & $87.96925804$ \\
Venus & $11.93672$ & $224.70082399$ \\
Erde & $111.28179$ & $365.25638350$ \\
Mars & $71.24145$ & $686.9796186$ \\
Jupiter & $124.67781$ & $4332.5963076$ \\
Saturn & $150.38010$ & $10758.9698400$ \\
Uranus & $197.54244$ & $30688.7126872$ \\
\end{tabular}} \end{center}

\begin{center} \Needspace{140pt}
{\small\setlength{\tabcolsep}{6pt}
\begin{tabular}{l|r|r|r}
 & \shortstack{Große\\Halbachse} & \shortstack{Exzentrizität\\1801} & \shortstack{Säkulare Änderung} \\
\hline
Merkur & $0.3870981$ & $0.20551494$ & $+ \text{c}.000003867$ \\
Venus & $0.7233323$ & $0.00685298$ & $- 0.000062711$ \\
Erde & $1.0000000$ & $0.01685318$ & $- 0.00\text{c}041632$ \\
Mars & $1.5236935$ & $0.09313400$ & $+ 0.000090176$ \\
Jupiter & $5\;2027911$ & $0.04817840$ & $+ 0.000159350$ \\
Saturn & $9.5387705$ & $0.05616830$ & $- 0.000312402$ \\
Uranus & $19.1833050$ & $0.04667030$ & $- 0.000025072$ \\
\end{tabular}} \end{center}

\SourceLeaf{507}
\begin{center} \Needspace{140pt}
{\small\setlength{\tabcolsep}{6pt}
\begin{tabular}{l|r|r}
 & \shortstack{Sonnennähe 1801} & \shortstack{Siderische\\Säkularbewegung} \\
\hline
Merkur & $82^{\circ}6256$ & $+ 1801''10$ \\
Venus & $142.9077$ & $- 826.63$ \\
Erde & $110.5571$ & $+ 3641.40$ \\
Mars & $369.3407$ & $+ 4884.05$ \\
Jupiter & $12.3812$ & $+ 2048.95$ \\
Saturn & $99.0549$ & $+ 5978.60$ \\
Uranus & $185.9574$ & $+ 738.69$ \\
\end{tabular}} \end{center}

\begin{center} \Needspace{140pt}
{\small\setlength{\tabcolsep}{6pt}
\begin{tabular}{l|r|r}
 & \shortstack{Neigung der Bahn\\1801} & \shortstack{Säkulare Änderung} \\
\hline
Merkur & $7^{\circ}78058$ & $+ 56''12$ \\
Venus & $3.76936$ & $- 14.05$ \\
Mars & $2.05663$ & $- 0.47$ \\
Jupiter & $1.46034$ & $- 69.78$ \\
Saturn & $2.77102$ & $- 47.88$ \\
Uranus & $0.85990$ & $+ 9.67$ \\
\end{tabular}} \end{center}

\begin{center} \Needspace{140pt}
{\small\setlength{\tabcolsep}{6pt}
\begin{tabular}{l|r|r}
 & \shortstack{Aufsteigender Knoten\\1801} & \shortstack{Siderische\\Säkularbewegung} \\
\hline
Merkur & $51^{\circ}0651$ & $- 2414''41$ \\
Venus & $83.1972$ & $- 5770.99$ \\
Mars & $53.3605$ & $- 7186.65$ \\
Jupiter & $109.3624$ & $- 4869.04$ \\
Saturn & $124.3662$ & $- 6995.25$ \\
Uranus & $10.9488$ & $- 11104.81$ \\
\end{tabular}} \end{center}

Von den hier angegebenen Elementen der vier neuen Planeten sind die der Ceres die 11., die der Pallas die 9. und die der Juno die 6. von \textsc{Gauss}. Für Vesta sind es diejenigen, die Herr \textsc{Burkhardt} in der Connoissance des tems für 1809 angegeben hat und für die erst wenige Beobachtungen benutzt worden waren. Wir lassen sie deshalb hier weg, da in Deutschland längst genauere bekannt sind (Monatl. Corresp. Febr. 1808; Gött. gel. Anz. 1808 St. 14. 40. 107).

Auch das \emph{dritte Buch} über die Bewegungsgesetze hat nur ein paar kleine Zusätze erhalten. Sie betreffen die Darstellung der ersten Grundsätze der Dynamik in Gleichungen, das Prinzip der kleinsten Wirkung, den Begriff \emph{Masse} und die beiden Arten des Gleichgewichts.

Im \emph{vierten Buch} über die Theorie der allgemeinen Schwere ist das Kapitel über die Störungen der elliptischen Bewegungen der Planeten zum zweiten geworden. Hinzugekommen sind zwei Kapitel, eines über die Trabanten des Saturn und Uranus, das andere über die Anziehung zwischen den kleinsten Teilen der Körper, die nur in unmerklichen Entfernungen merklich ist (attraction moléculaire). Dieses letzte Kapitel gibt einen Überblick über die Theorie, die \textsc{Laplace} in zwei Ergänzungen zu seiner Mécanique céleste (De l'action capillaire und Supplément à la théorie de l'action capillaire) entwickelt hat. Er hat sie mit so überraschend glücklichem Erfolg zur mathematischen Erklärung der Lichtbrechung, der Erscheinungen in Kapillarröhren, der Anziehung und Abstoßung kleiner auf einer Flüssigkeit schwimmender Körper, der Haftung einer Scheibe an einer Flüssigkeit, des Schwimmens kleiner fester Körper in einer spezifisch leichteren Flüssigkeit und so weiter angewandt. Diese Entdeckungen eröffnen eine neue Epoche in der Physik. Sie bahnen den Weg zu einer Wissenschaft, die für die Natur in Bezug auf die kleinsten Teile der Körper das sein wird, was die Gravitationstheorie für die Natur im Großen ist. Eine ausführlichere Besprechung davon müssen wir uns aber für eine andere Gelegenheit aufsparen.

\SourceLeaf{508}
Von den kleineren Änderungen und Zusätzen in diesem Abschnitt nennen wir noch die folgenden. Die große Gleichung des Saturn steigt auf $9111''41$ und ihre Periode beträgt $921\frac{1}{2}$ Jahre. Die große Gleichung des Jupiter beträgt $3720''36$. Die Massen der Planeten verhalten sich zur Masse der Sonne wie 1 zu 2025810 bei Merkur, zu 356632 bei Venus, zu 337086 bei der Erde, zu 2546320 bei Mars, zu 1067.09 bei Jupiter, zu 3534.08 bei Saturn und zu 19504 bei Uranus. Der scheinbare Durchmesser des Jupiter in der Entfernung 1 beträgt $599''151$. Die Schwere eines Körpers am Erdäquator, am Jupiteräquator und am Sonnenäquator verhält sich wie die Zahlen 1000, 2566, 27933. In die Theorie des Mondes sind mehrere Zusätze eingefügt. Sie betreffen die Ungleichheiten, die von der Abplattung der Erde und von der Parallaxe der Sonne abhängen, ferner die säkularen Ungleichheiten in der Bewegung des Knotens und der Apsiden und die vor einigen Jahren entdeckte Gleichung mit einer Periode von 184 Jahren. --- Auch der Abschnitt über die Jupitertrabanten ist an mehreren Stellen ausführlicher geworden.

Das \emph{fünfte Buch} enthält die Geschichte der Astronomie. Auch hier sind es nur Umrisse, doch von solcher Hand gezeichnet, sind sie für Kenner wie für Liebhaber von großem Interesse. „Eine Darstellung der Fortschritte der erhabensten Naturwissenschaft“, sagt \textsc{Laplace}, „von Fortschritten, die selbst inmitten der Umwälzungen der Reiche stetig weitergehen, mag über die Unglücksfälle hinwegtrösten, deren Schilderungen die Annalen aller Völker füllen.“ Auch dieser Teil des Werks ist um manche Zusätze und interessante Überlegungen bereichert. Dazu gehören die Bemerkungen über die Spuren astronomischer Kenntnisse, die man bei den indigenen Völkern Mexikos und Perus vorfand. Erstere kannten die Länge des tropischen Jahres sehr genau. Sie verwendeten eine Schaltmethode, die dafür $365\frac{25}{104}$ Tage oder 365 Tage, 5 Stunden, 46 Minuten und 9 Sekunden voraussetzt. Dabei ist sehr bemerkenswert, dass ihnen die Einteilung der Zeit in Wochen, die man bei allen Völkern der Alten Welt findet, unbekannt war und dass sie stattdessen einen Zeitraum von fünf Tagen hatten.

Am Schluss des Werks stehen noch sechs historische Anmerkungen. Die erste bezieht sich auf einige chinesische Beobachtungen von \textsc{Tcheou-Kong}, nach denen ungefähr im Jahr 1100 vor unserer Zeitrechnung die Schiefe der Ekliptik $26^{\circ}5563$ und die Rektaszension des Sterns $\varepsilon$ im Wassermann $297^{\circ}8096$ betrugen. Nach den Formeln von \textsc{Laplace} sollten sie für dieses Jahr $26^{\circ}5161$ beziehungsweise $298^{\circ}7265$ betragen. Um die letzteren Zahlen vollständig in Übereinstimmung zu bringen, müsste man nur noch 54 Jahre weiter zurückgehen. Die zweite betrifft einen Bericht, den uns \textsc{Geminus} über die Kenntnisse der Chaldäer von der Bewegung des Mondes überliefert hat. Die dritte Anmerkung betrifft die Beobachtung der Sonnenwende in Marseille durch \textsc{Pytheas}. Die vierte vergleicht die Angaben von \textsc{Hipparch} zur Bewegung des Mondes in Bezug auf die Sonne, die Erdnähe und den Knoten mit den Bestimmungen, die aus der Theorie der säkularen Ungleichheiten von \textsc{Laplace} folgen. Die fünfte vergleicht die größte Mittelpunktsgleichung der Sonne, die Lage der Sonnenferne, die Länge des Jahres und die Schiefe der Ekliptik nach den Bestimmungen der Araber mit den neuesten Angaben. Die sechste schließlich stellt die Bestimmungen der Schiefe der Ekliptik von \textsc{Tcheou-Kong} aus dem Jahr 1100 vor Christus, von \textsc{Pytheas} aus dem Jahr 350 vor Christus, von \textsc{Ibn Junis} aus dem Jahr 1000, von \textsc{Cocheou-Kong} aus dem Jahr 1280 und von \textsc{Ulug-Beigh} aus dem Jahr 1437 sowie die neuesten von 1801 den Ergebnissen der Theorie gegenüber.
